\documentclass[11pt]{article}
\usepackage[a4paper,margin=26mm]{geometry}
\usepackage{amsmath,amssymb,amsthm,mathtools,booktabs,hyperref,xurl}
\setlength{\emergencystretch}{3em}
\hypersetup{colorlinks=true,linkcolor=blue,citecolor=blue,urlcolor=blue}
\newtheorem{theorem}{Theorem}[section]
\newtheorem{lemma}[theorem]{Lemma}
\newtheorem{corollary}[theorem]{Corollary}
\newcommand{\CB}{\mathrm{CB}}
\newcommand{\DB}{\mathrm{DB}}
\newcommand{\WR}{\mathrm{WR}}
\newcommand{\RR}{\mathbb R}
\newcommand{\NN}{\mathbb N}
\title{Sharp molecularity bounds for balanced mass-action realizations on a two-species conservation class}
\author{Research manuscript: human authorship and responsibility pending confirmation}
\date{November 26, 2026 --- v0.4}
\begin{document}
\maketitle
\begin{abstract}
We give an exact molecularity criterion for deterministic conservative mass-action realizations on a prescribed two-species population simplex. For a fixed total $N>0$, an interior balancing point $a\in(0,N)$, and a molecularity cap $D$, a scalar field $f$ admits a complex-balanced realization of molecularity at most $D$ if and only if it admits a homogeneous detailed-balanced realization of degree $D$, equivalently if $f(t)=(a-t)\sum_{i=0}^{D-1}c_i t^i(N-t)^{D-1-i}$ with $c_i\geq0$. Necessity follows from nonnegative prefix flux sums; sufficiency gives explicit adjacent reversible reaction pairs. This yields a sharp unbounded gap between polynomial degree and the minimum molecularity of a balanced realization: a family of cubic fields on the unit simplex with a unique simple attracting interior equilibrium requires arbitrarily high molecularity. We derive exact even/odd degree thresholds, including equality. Motivated by open problems in mathematical epidemiology, we also establish the distinction between global weakly reversible realization and realization on a fixed population class for the standard SIS, SIR and SIRS models. The key mathematical statements are formalized in Lean 4 with Mathlib, with independent rational, symbolic and linear-programming checks. We distinguish deterministic restricted-domain equality from global field equality, reaction-graph equality, stochastic equality and physical elementary-reaction feasibility. The manuscript is AI-assisted and has not been peer reviewed.
\end{abstract}

\section{Introduction}
Reaction-network realization asks whether a polynomial vector field can be generated by a mass-action network having a specified structural or balance property. The network is not determined uniquely by its vector field. Global realization methods for reversible, weakly reversible, detailed-balanced and complex-balanced systems have been extensively developed, including linear-programming and flux formulations \cite{szed,efficient,realizations,boros}. Disguised toric systems on collinear complexes, including adjacent-pair constructions and the distinction between stable and globally disguised toric systems, were studied in \cite{disguised}. We do not claim novelty for the general realization problem, collinear networks, or the use of balanced fluxes. Endotacticity is a known necessary condition in the global weakly reversible realization setting \cite{kothari}.

Here the domain of equality is instead a \emph{single prescribed} invariant population class. Restricting to $S+I=N$ makes homogeneous monomials of different degrees dependent: multiplication by $(S+I)/N$ changes a polynomial globally without changing its restriction. This flexibility prompts a quantitative question: how much molecularity is necessary to realize a simple scalar field with balance? In particular, does the degree of the scalar polynomial impose a uniform upper bound on the molecularity needed? Our exact finite-degree criterion and explicit cubic family answer this question negatively. A second motivation is the weak reversibility and weakly reversible deficiency-zero realization problems raised for epidemic models in \cite{avram24}; we give precise same-species results for classical closed models, rather than assert a solution of the unrestricted open problems.

\section{Definitions and scope}
A finite ordinary deterministic mass-action network consists of species concentrations $x\in\RR_{>0}^q$, integer complexes $y\in\NN^q$, directed reactions $y\to z$ and positive constant rates $k_{yz}$. Its field is $F(x)=\sum_{y\to z}k_{yz}x^y(z-y)$. A zero rate in a matrix representation denotes an absent reaction; diagonal terms can be omitted. A network is weakly reversible if every edge has a return path. It is complex balanced (CB) at a positive state if the total incoming and outgoing fluxes at every complex agree. It is detailed balanced (DB) there if the fluxes of every opposing pair agree. These are properties of a realization, not necessarily of every network giving the same field.

For the main theorem we restrict to two species $S,I$, with \emph{each reaction} preserving the unweighted total $S+I$. The maximum molecularity is the maximum sum of stoichiometric coefficients in a complex participating in a reaction; for such a conservative network the source and target have the same degree. An arbitrary network with maximum molecularity at most $D$ is thus exactly represented by nonnegative rate matrices $K^{(d)}=(k_{ij}^{(d)})_{0\leq i,j\leq d}$, $0\leq d\leq D$, on the complexes $Y_i^{(d)}=(d-i)S+iI$. Missing complexes have zero incident rates. Adding or dropping zero-rate entries or self-loops does not change the field or balance. In particular, there is no assumption of a given reaction graph or of nearest-neighbor reactions in the necessity argument.

Fix $N>0$, $0<a<N$. Define $B_{d,i}(t)=t^i(N-t)^{d-i}$ for $0\leq i\leq d$. On the positive simplex, writing $t=I$ and $S=N-t$, the infected derivative from degree $d$ is
\begin{equation}\label{field}
F_d(t)=\sum_{i,j=0}^{d} k^{(d)}_{ij}(j-i)B_{d,i}(t),
\qquad F(t)=\sum_{d=0}^D F_d(t).
\end{equation}
The susceptible derivative is $-F(t)$ because each reaction conserves the total. Consequently scalar equality here is equality of the complete two-species vector field on this simplex, not merely of its equilibrium or of one trajectory. Balance at $(N-a,a)$ is a condition on the same constant rate matrices. Rates may be chosen using the prescribed $N$ and $a$; a different total is not required to have the same restricted field.

\section{Exact bounded-molecularity criterion}
\begin{theorem}[Balanced realizations at a prescribed total and equilibrium]\label{main}
Let $N>0$, $0<a<N$ and $D\geq0$ be an integer. For any function $f:(0,N)\to\RR$, the following statements are equivalent:
\begin{enumerate}
\item The field $(-f,f)$ on $S+I=N$ has a finite conservative CB realization at $(N-a,a)$ of maximum molecularity at most $D$.
\item It has a conservative DB realization at that point with every participating complex of molecularity $D$.
\item There exist $c_0,\ldots,c_{D-1}\geq0$ such that
\begin{equation}\label{factor}
f(t)=(a-t)\sum_{i=0}^{D-1}c_i B_{D-1,i}(t),\qquad 0<t<N.
\end{equation}
\end{enumerate}
For $D=0$ the sum is empty and the field is zero. The empty network is allowed for the zero field.
\end{theorem}
\begin{proof}
For a degree $d$ matrix, put $w_{ij}=k^{(d)}_{ij} B_{d,i}(a)$ and $v_i=\sum_j w_{ij}(j-i)$. Complex balance means $\sum_jw_{ij}=\sum_jw_{ji}$ for every $i$. For every function $g$ on the index set, rearranging the finite sums gives
\begin{equation}\label{potential}
\sum_{i,j} w_{ij}(g(j)-g(i))=0.
\end{equation}
Taking $g(i)=i$ yields $\sum_i v_i=0$. For an integer cut $k$, take $g(i)=\min(i,k)$. The pointwise inequality
\[
\min(j,k)-\min(i,k)\leq \mathbf 1_{i<k}(j-i)
\]
holds whether $i<k$ or $i\geq k$. Multiplication by the nonnegative fluxes and summation, using \eqref{potential}, gives
\begin{equation}\label{prefix}
P_k:=\sum_{i<k}v_i\geq0\quad (0\leq k\leq d+1).
\end{equation}
Let $b_i(t)=B_{d,i}(t)/B_{d,i}(a)$. Since $P_0=P_{d+1}=0$ and $v_i=P_{i+1}-P_i$, summation by parts in \eqref{field} gives
\[
F_d(t)=\sum_i v_i b_i(t)=\sum_{i=0}^{d-1}P_{i+1}(b_i(t)-b_{i+1}(t)).
\]
Direct cancellation of the positive denominators shows
\[
b_i(t)-b_{i+1}(t)
=(a-t)\frac{N B_{d-1,i}(t)}{a^{i+1}(N-a)^{d-i}}.
\]
Thus each degree contributes a factor of form \eqref{factor} with $D=d$ and nonnegative coefficients. The $d=0$ contribution vanishes. For $d\leq D$, elevate each factor using
\begin{equation}\label{elevate}
 B_{m,i}(t)=\frac{B_{m+1,i}(t)+B_{m+1,i+1}(t)}{N}.
\end{equation}
The degree increases by one, the restricted function stays unchanged, and nonnegative coefficients remain nonnegative. Adding the elevated factors proves (1)$\Rightarrow$(3).

Conversely, for each $i$ with $c_i>0$ use the pair
\begin{equation}\label{pair}
Y_i^{(D)}\ \xrightleftharpoons[c_i(1-a/N)]{c_i a/N}\ Y_{i+1}^{(D)}.
\end{equation}
Both rates are positive; the two complexes have molecularity $D$. Their equilibrium fluxes coincide since $(a/N)(N-a)=(1-a/N)a$. The pair's contribution is
\[
 c_i\big[(a/N)(N-t)-(1-a/N)t\big] B_{D-1,i}(t)
 =(a-t)c_iB_{D-1,i}(t).
\]
Summing over all positive coefficients realizes the field. Omit pairs with zero coefficient. Thus (3)$\Rightarrow$(2). Detailed balance implies complex balance, proving (2)$\Rightarrow$(1).
\end{proof}

\begin{corollary}[Unique coefficient certificate]\label{coeff}
For $D\geq1$, the coefficients in \eqref{factor}, when they exist, are unique. Hence for a prescribed polynomial $f$ and $a$, feasibility at a cap $D$ is exactly the polynomial equality test and nonnegativity of these coefficients.
\end{corollary}
\begin{proof}
Cancel $a-t$ on $(a,N)$. Divide the two candidate expansions by $(N-t)^{D-1}$, and set $x=t/(N-t)$. Equality becomes equality of two ordinary polynomials in $x$ on an infinite interval. Their coefficients coincide. No division is needed at $t=a$, where both fields vanish.
\end{proof}

The certificate has $D$ nonnegative scalar entries, whereas a direct search for general rate matrices up to $D$ introduces $\sum_{d=0}^D(d+1)^2$ entries. Equation \eqref{pair} constructs a realization with at most $2D$ nonzero directed reactions; no claim is made here that this is the smallest possible number of reactions. The feasibility sets are nested under degree elevation. Every nonzero feasible field has the sign of $a-t$ in the entire open simplex and a strictly negative derivative at $a$ if it is differentiable there, because the factor sum is strictly positive on $(0,N)$. This is consistent with classical balance theory; the new quantitative question is the smallest admissible $D$.

\section{Cubic fields with unbounded required molecularity}
Consider the unit simplex and the cubic family
\begin{equation}\label{cubic}
f_{\varepsilon,a}(t)=(a-t)q_\varepsilon(t),\qquad q_\varepsilon(t)=(t-\tfrac12)^2+\varepsilon,
\end{equation}
where $0<a<1$ and $\varepsilon>0$. Its only interior zero is $a$; the sign is inward towards $a$ and $f'_{\varepsilon,a}(a)=-q_\varepsilon(a)<0$. The polynomial degree is exactly three, independent of $\varepsilon$.
\begin{theorem}[Sharp parity-dependent threshold]\label{sharp}
For an integer cap $D\geq3$, set $n=D-1$. The equivalent conditions of Theorem~\ref{main} hold for $f_{\varepsilon,a}$ if and only if
\begin{equation}\label{threshold}
\varepsilon\geq \begin{cases}1/[4(n-1)]& n\text{ even},\\1/(4n)&n\text{ odd}.\end{cases}
\end{equation}
The equivalence includes equality, with zero-rate pairs omitted. Caps $D\leq2$ cannot realize this cubic.
\end{theorem}
\begin{proof}
Use the ordinary normalized Bernstein polynomials $\beta_{n,i}(t)=\binom ni t^i(1-t)^{n-i}$. The binomial theorem gives $\sum_i\beta_{n,i}=1$, $\sum_i i\beta_{n,i}=nt$, and $\sum_i i(i-1)\beta_{n,i}=n(n-1)t^2$. Expanding the quadratic using these identities yields coefficients
\begin{align*}
b_{n,i}&=\varepsilon+\tfrac14-\frac{i}{n}+\frac{i(i-1)}{n(n-1)}\\
&=\varepsilon+\frac{(2i-n)^2-n}{4n(n-1)}.
\end{align*}
The unnormalized coefficients in \eqref{factor} are $c_i=\binom ni b_{n,i}$, whose signs are those of $b_{n,i}$. By Corollary~\ref{coeff}, no alternate coefficient vector at the same degree can avoid their signs. For even $n$, the minimum of $(2i-n)^2$ is zero, attained at $i=n/2$. For odd $n$ it is one, attained at the two central integers. Substitution proves \eqref{threshold}. A field from matrices of degree at most two has degree at most two after restriction, so cannot equal a cubic on the interval.
\end{proof}
\begin{corollary}[Unbounded molecularity gap]\label{gap}
There is no function of the polynomial degree alone that bounds the minimum molecularity required for conservative complex-balanced or detailed-balanced realizations on the unit simplex, even for polynomial fields having one simple attracting interior equilibrium. In fact the family \eqref{cubic} has finite DB realizations for every $\varepsilon>0$, but their minimum possible molecularity diverges as $\varepsilon\downarrow0$.
\end{corollary}
\begin{proof}
The thresholds in \eqref{threshold} tend to zero as $n\to\infty$, and are strictly positive at every fixed $n$. Given any finite cap $D$, choose a higher cap $2m+1\geq\max(3,D)$ and $0<\varepsilon<1/[4(2m-1)]$. Infeasibility at $2m+1$ implies infeasibility at $D$. Choosing a still larger degree gives feasibility and the explicit DB construction. More precisely for small $\varepsilon$ the minimum degree differs from $1/(4\varepsilon)$ by a bounded additive amount (the parity and endpoint rounding follow from \eqref{threshold}).
\end{proof}

\begin{table}[ht]
\centering
\caption{Exact minimum cap for the cubic family on $N=1$. The location $a$ does not enter the threshold.}
\begin{tabular}{rrr}
\toprule
$\varepsilon$ & minimum $D$ & polynomial degree\\
\midrule
$1/4$ & 3 & 3\\
$1/12$ & 4 & 3\\
$1/20$ & 6 & 3\\
$1/100$ & 26 & 3\\
$1/1000$ & 252 & 3\\
\bottomrule
\end{tabular}
\end{table}

\section{Global versus fixed-class equality}\label{sec:global}
For clarity, an extension of \eqref{cubic} to the full positive quadrant is the homogeneous cubic
\begin{equation}\label{extension}
G(S,I)=(a(S+I)-I)\big[(I-(S+I)/2)^2+\varepsilon(S+I)^2\big].
\end{equation}
It restricts to \eqref{cubic} on $S+I=1$. At scale $r>0$, it equals $r^3f_{\varepsilon,a}(t)$ at $(r(1-t),rt)$. In contrast to the fixed-class problem, rates in a global realization cannot depend on $r$.
\begin{theorem}[Homogeneous extension separation]\label{sep}
For $0<\varepsilon<1/4$ and $0<a<1$, the cubic \eqref{extension} cannot have any finite conservative CB realization on the entire positive quadrant balanced at $(1-a,a)$, regardless of molecularity cap. Yet its restriction to $S+I=1$ does have a finite conservative DB realization. The same obstruction applies to an arbitrary positive balancing point: scaling it to total one preserves balance within each degree, and the restricted field has the unique positive equilibrium $a$.
\end{theorem}
\begin{proof}
For every matrix of degree $d$, the field contribution obeys $F_d(r(1-t),rt)=r^d F_d(1-t,t)$. Suppose a finite conservative global CB realization exists. Equality with \eqref{extension} for all $r>0$ and fixed $0<t<1$ gives an identity of polynomials in $r$. The coefficient of $r^3$ must be $f_{\varepsilon,a}(t)$ (if degree three is absent, that coefficient is zero, contradicting the nonzero field). Complex balance separates across degrees because each reaction preserves total, so the degree-three matrix is itself balanced. It would therefore be a degree-three CB realization on the unit simplex. Theorem~\ref{sharp} at $n=2$ rules this out. Existence on the unit simplex follows from Corollary~\ref{gap}. Finally scaling a balanced point by a positive constant multiplies every flux at a vertex of degree $d$ by the same power of that constant; thus incoming and outgoing totals still agree. A balanced point must have zero field, as follows from \eqref{potential} with $g(i)=i$, so the unit-total point must be $a$.
\end{proof}
This illustrates why global disguised-toric results cannot simply be transferred to a single conservation class. A fixed-class representation may use extra degrees precisely because those degrees become dependent only after restriction.

\section{Epidemiological applications}\label{sec:epi}
The following statements give a complementary, weaker structural condition (weak reversibility, not balance). We use finite arbitrary integer complexes and positive rates; the competing network need not be conservative outside the prescribed domain. The extremal-source argument is a directional form of known endotactic geometry \cite{classes,kothari}; it is included to make the applications self-contained.
\begin{lemma}[Directional obstruction]\label{dir}
Let $J$ be a subset of species, $h(y)=\sum_{j\in J}y_j$, and $F_J=\sum_{j\in J}F_j$. In a weakly reversible network with a height-changing edge, the height-changing edges at minimum active source height increase height, and those at maximum active source height decrease height. Along $x_j(t)=a_jt$ for $j\in J$ and $x_i(t)=b_i(t)$ otherwise, with $a_j>0$ and $b_i$ continuous and positive at zero, $F_J$ is strictly positive for sufficiently small positive $t$. If the $b_i$ are constant positive, it is strictly negative for sufficiently large $t$. With no height-changing edge the projection vanishes.
\end{lemma}
\begin{proof}
An edge from a minimum active height cannot go below it: its return path would contain a height-increasing edge from a lower active source. The maximum statement is the reversed argument. Each contribution to $F_J$ equals $t^{h(y)}k_{yz}(h(z)-h(y))\prod_{j\in J}a_j^{y_j}\prod_{i\notin J}b_i(t)^{y_i}$. Divide by the smallest active power of $t$ and take the limit at zero; the leading coefficient is a nonempty sum of positive terms. For constant $b_i$, divide instead by the greatest active power and take the limit at infinity, where the leading coefficient is strictly negative.
\end{proof}
\begin{theorem}[Exact SIS threshold]\label{sis}
For $N,\beta,\gamma>0$, the closed SIS field $f(t)=\beta(N-t)t-\gamma t$ admits a same-species weakly reversible realization on the entire positive simplex $S+I=N$ if and only if $\beta N>\gamma$. When this holds, an explicit conservative reversible network is
\[
S+I\ \xrightleftharpoons[\gamma/N]{\beta-\gamma/N}\ 2I.
\]
The corresponding global field $\beta SI-\gamma I$ cannot have a same-species weakly reversible realization throughout the positive quadrant.
\end{theorem}
\begin{proof}
If $\beta N\leq\gamma$, then $f(t)<0$ for all $0<t<N$, contradicting Lemma~\ref{dir} with $I=t$, $S=N-t$ and $J=\{I\}$. If $\beta N>\gamma$, the two displayed rates are strictly positive. The network is weakly reversible, and its infected derivative equals $(\beta-\gamma/N)SI-(\gamma/N)I^2=\beta SI-\gamma I$ on the simplex. Conservation gives the equality of the other component as well. For global nonexistence, choose $S=(\gamma+1)/\beta$ and $I=t$. The target infected derivative is exactly $t$, which cannot equal the identically zero or eventually negative projection of a weakly reversible network. The reversible pair also balances at $a=N-\gamma/\beta$; its two complexes and one linkage class span a one-dimensional reaction space, giving deficiency zero by the standard count.
\end{proof}
\begin{theorem}[SIR/SIRS obstruction]\label{sir}
For standard closed SIR and SIRS fields with positive transmission and recovery rates, no finite same-species weakly reversible mass-action network agrees with the field on the entire positive population simplex. This remains true with any nonnegative immunity-loss rate.
\end{theorem}
\begin{proof}
The infected derivative is $I(\beta S-\gamma)$. Choose $S_0=\min(N/2,\gamma/(2\beta))>0$. Then $S_0<N$ and $\beta S_0<\gamma$. Along $I=t$, $S=S_0$, $R=N-S_0-t$, for $0<t<N-S_0$, the infected derivative is strictly negative. The other species have positive limiting values at zero, so Lemma~\ref{dir} forbids this. The immunity-loss term does not change the infected derivative.
\end{proof}

The broader genuine-infection reaction class in \cite{avram24} suggests an additional general projection argument: scaling the infected group by $t$ while holding other concentrations fixed can produce an affine infected derivative with positive slope, conflicting with the maximum-source sign. The Lean development verifies the full mass-action affine-projection obstruction, but the reduction of every syntactic subclass of the published epidemic model definition to that premise is not claimed as a machine-checked solution of Open Problems 2 or 5. In particular, transfer-only degeneracies, additional species, alternate kinetic laws and equivalence only along selected trajectories are not included in the stated applications.

\section{Verification, prior art, and limitations}
\paragraph{Formal verification.} The project uses Lean 4.28.0 and Mathlib v4.28.0. The file \texttt{HomogeneousCuts\_v04.lean} (51 theorem declarations) contains the finite degree-indexed rate matrices, flux balance, restricted fields and realizability predicates. Theorems \ref{main} and \ref{sharp}, uniqueness of the coefficients, degree elevation, infeasibility for every cap $D\leq2$, existence of a finite DB realization for every $\varepsilon>0$, and global-versus-fixed-class separation at each prescribed $0<a<1$ have corresponding kernel-checked statements. For example, the declarations \texttt{mixed\_cb\_db\_exact}, \texttt{mixed\_exact}, and \texttt{cubic\_finite\_db} establish the principal existence equivalences. The unbounded-gap statement is also checked directly for $a=1/3$; the general parity thresholds and finite existence quantify over arbitrary interior $a$. The extension to an arbitrary positive global balancing point follows from the elementary scaling argument in Theorem~\ref{sep}; this normalization argument is not a separate machine-checked declaration. \texttt{EpidemicApplications.lean} directly quantifies arbitrary finite reaction sets with return paths and positive rate constants; the exact SIS threshold, global SIS obstruction and SIR/SIRS obstruction are kernel-checked. The reversible pair's deficiency count is an analytic statement, not a Lean deficiency-library calculation. All three standalone source files compile with exit code zero. There are 90 distinct theorem declarations in the two final developments (the older 31-theorem file is contained in the 39-theorem epidemic file and is preserved for provenance, not counted twice). The logs and SHA-256 hashes accompany the source files; the audited declarations depend only on standard Lean logical axioms. Kernel checking concerns the mathematical statements under the model assumptions, not biological reality.

\paragraph{Independent verification.} A separate Python program checks exact rational flux and adjacent-pair identities on 800 generated balanced mixed networks at three sample points each (2,400 identities); it compares the criterion against 700 independent full directed equilibrium-flux linear-program feasibility tests at degrees one through seven. The program also checks 297 exact coefficient-threshold cases at degrees two through 100 and 11 symbolic quadratic expansions. All recorded tests pass. Finite tests support implementation fidelity; the universal claims rest on proofs, not numerical sampling or LP solver precision.

\paragraph{Novelty assessment.} Linear-program realizability, zero-deficiency algorithms, classical Bernstein degree elevation and polynomial positivity are established mathematics \cite{szed,efficient,algorithm,buxton}. The collinear-complex examples in \cite{disguised}, especially Sections 4, 5 and 7, and the more general global flux-cone framework in \cite{boros} are closely related. The single-degree collinear realization mechanism overlaps substantially with these works, and is not claimed as a standalone priority result. The contribution under review is specifically the mixed-degree bounded-molecularity certificate on one fixed simplex, the sharp cubic degree thresholds and unbounded molecularity gap, and the associated distinction from global conservative realizability. We do not claim that the general phenomenon of stable but globally non-disguised-toric fields is new. The epidemic observations are applications of known directional ideas; their inclusion is not a claim of priority over all previous epidemiological realization results. See the accompanying search and comparison record for coverage and limitations. A literature search cannot establish absolute priority, and independent expert assessment remains necessary.

\paragraph{Physical and mathematical limitations.} The conservation law is the unweighted sum of two modeled concentrations, not arbitrary chemical elemental conservation. Only ordinary mass-action kinetics on the same species are considered. High molecularity is a mathematical representation cost, not a prediction of simultaneous many-body elementary collisions. Hidden intermediates, timescale reductions, generalized kinetics and stochastic networks may lead to different conclusions. No stochastic equality is claimed. No structural graph equivalence between two realizations is asserted. The previous exploratory draft also considered minimum reaction counts and an unbounded-degree classification for all inward polynomials, which are not asserted as new fully formalized theorems in this focused version. That draft and its remaining proof tasks are retained separately.

\paragraph{Research integrity.} AI contributed substantially to literature discovery, derivations, proof coding, numerical verification and manuscript preparation. No human authorship or institutional affiliation is asserted before confirmation of actual contribution and responsibility. No peer review or journal acceptance is claimed. This is a research manuscript and reproducibility package, not a submitted paper.

\begin{thebibliography}{12}
\bibitem{avram24} F. Avram, R. Adenane and M. Neagu, Advancing mathematical epidemiology and chemical reaction network theory via synergies between them. \emph{Entropy} 26 (2024), 936. \url{https://doi.org/10.3390/e26110936}.
\bibitem{szed} G. Szederk\'enyi and K. M. Hangos, Finding complex balanced and detailed balanced realizations of chemical reaction networks. \emph{J. Math. Chem.} 49 (2011), 1163--1179. \url{https://doi.org/10.1007/s10910-011-9804-9}.
\bibitem{efficient} G. Craciun, J. Jin and P. Y. Yu, An efficient characterization of complex-balanced, detailed-balanced, and weakly reversible systems. \emph{SIAM J. Appl. Math.} 80 (2020), 183--205. \url{https://doi.org/10.1137/19M1244494}.
\bibitem{realizations} G. Craciun, M. D. Johnston, G. Szederk\'enyi, E. Tonello, J. T\'oth and P. Y. Yu, Realizations of kinetic differential equations. \emph{Math. Biosci. Eng.} 17 (2020), 862--892. \url{https://aimspress.com/article/id/4373}.
\bibitem{disguised} L. Brustenga i Moncus\'i, G. Craciun and M.-S. Sorea, Disguised toric dynamical systems. \emph{J. Pure Appl. Algebra} 226 (2022), 107035. \url{https://doi.org/10.1016/j.jpaa.2022.107035}.
\bibitem{boros} B. Boros, G. Craciun, O. Henriksson, J. Jin and D. Rojas La Luz, A flux-based approach for analyzing the disguised toric locus of reaction networks. \emph{Adv. Appl. Math.} 182 (2027), 103160, available online 2026. \url{https://doi.org/10.1016/j.aam.2026.103160}; arXiv:2510.03621.
\bibitem{algorithm} G. Craciun, J. Jin and P. Y. Yu, An algorithm for finding weakly reversible deficiency zero realizations of polynomial dynamical systems. \emph{SIAM J. Appl. Math.} 83 (2023), 1717--1737. \url{https://doi.org/10.1137/22M1499558}.
\bibitem{buxton} N. Buxton, G. Craciun, A. Deshpande and C. Pantea, Weakly reversible deficiency zero realizations of reaction networks. \emph{J. Nonlinear Sci.} 35 (2025), 104. \url{https://doi.org/10.1007/s00332-025-10200-4}.
\bibitem{classes} D. F. Anderson, J. D. Brunner, G. Craciun and M. D. Johnston, On classes of reaction networks and their associated polynomial dynamical systems (2020). \url{https://people.math.wisc.edu/~dfanderson/papers/2020/ABCJ_Endo.pdf}.
\bibitem{kothari} A. Deshpande, J. Jin and S. Kothari, Realizations Through Weakly Reversible Networks and the Globally Attracting Locus. \emph{SIAM J. Appl. Dyn. Syst.} 25 (2026), 42--76. \url{https://doi.org/10.1137/25M1745027}; arXiv:2409.04802.
\end{thebibliography}
\end{document}
